(1) 温度均匀为 \(T\) 的无限大平板，厚度 \(L\)，灰体吸收系数 \(\kappa>0\) 为常数且无散射。其频率积分后的局域热平衡源函数（每立体角辐亮度）为 \(B=\sigma_{\rm SB}T^4/\pi\)。两侧均为真空且无入射辐射，令 \(\mu=\cos\vartheta\)、\(z\in[0,L]\)。从 \(\mu\,dI/dz=\kappa(B-I)\) 推导向上与向下的定常辐亮度随位置的精确表达式。
(2) 令总光学厚度 \(\tau_0=\kappa L\)，定义 \(E_n(x)=\int_1^\infty e^{-xt}t^{-n}\,dt\)。从半球角积分求每个表面的向外辐射通量 \(F_{\rm face}\)，给出 \(\tau_0\ll1\) 与 \(\tau_0\gg1\) 的渐近，并解释薄板在无散射时为何不立即呈完全黑体。
(3) 定义角平均 \(J(z)=(4\pi)^{-1}\int I(z,\Omega)\,d\Omega\)。求 \(J(z)\)，从输运方程导出局部辐射净冷却率 \(q(z)=4\pi\kappa[B-J(z)]\)，证明其为非负并关于板中面对称。
(4) 不用逐角穿透率而直接积分上一问的体冷却率，证明其等于两表面的总能量出流率 \(2F_{\rm face}\)。核验 \(\tau_0\to0\) 的体功率标度、无限光学深度极限及灰体局域热平衡模型的适用条件。
